\documentclass[conference]{IEEEtran}
\IEEEoverridecommandlockouts

\usepackage[T1]{fontenc}
\usepackage[utf8]{inputenc}
\usepackage[ngerman]{babel}
\usepackage{amsmath}
\usepackage{booktabs}
\usepackage{array}
\usepackage{graphicx}
\usepackage{subcaption}
\usepackage{listings}
\usepackage{xcolor}
\usepackage{url}
\usepackage{hyperref}
\hypersetup{hidelinks}

% --- Listing style for shell / python ---------------------------------------
\definecolor{lstbg}{gray}{0.96}
\definecolor{lstkw}{rgb}{0.10,0.10,0.55}
\definecolor{lstcm}{rgb}{0.30,0.45,0.30}
\lstset{
  basicstyle=\ttfamily\scriptsize,
  backgroundcolor=\color{lstbg},
  keywordstyle=\color{lstkw}\bfseries,
  commentstyle=\color{lstcm}\itshape,
  breaklines=true,
  breakatwhitespace=true,
  columns=fullflexible,
  keepspaces=true,
  frame=single,
  framesep=3pt,
  rulecolor=\color{gray!40},
  xleftmargin=2pt,
  showstringspaces=false,
  postbreak=\mbox{\textcolor{gray}{$\hookrightarrow$}\space},
}
\lstdefinestyle{sh}{language=bash}
\lstdefinestyle{py}{language=Python}

\newcommand{\param}[1]{\texttt{#1}}

\begin{document}

\title{Reproduzierbare Aufbereitung häuslicher Blutdruck-Zeitreihen:\\
Methodik und Bedienung des Generators \texttt{generate\_bp\_tikz.py}}

\author{%
\IEEEauthorblockN{Erwin Grabler}
\IEEEauthorblockA{ARTCOM Vienna, Wien, Österreich\\
\texttt{OE1EKG}}}

\maketitle

\begin{abstract}
Häusliche Blutdruck-Selbstmessungen (HBPM) entstehen typischerweise mit
unregelmäßiger Messhäufigkeit: an manchen Tagen mehrere Messungen, an anderen
keine. Eine einfache Mittelung über alle Einzelmessungen gewichtet daher
mess\-intensive Tage über und verzerrt das Gesamtbild. Dieses Dokument
beschreibt die Methodik und die Bedienung von \texttt{generate\_bp\_tikz.py},
einem Kommandozeilenwerkzeug, das aus einer einfachen CSV-Datei zwei
aufeinander abgestimmte LaTeX/PGFPlots-Diagramme sowie maschinenlesbare
Statistik-CSVs erzeugt. Kernprinzip ist eine zweistufige Aggregation: zuerst
pro Kalendertag, dann optional in Blöcken fester Länge. Es werden die
Berechnungsvarianten (nach Kalendertagen gewichteter Mittelwert gegenüber
Median der Tagesmediane), die Ausreißerlogik sowie die erzeugten Datenformate
erläutert und sämtliche Parameter tabellarisch dokumentiert. Ziel ist eine
reproduzierbare, transparente und klinisch interpretierbare Darstellung, die
ausdrücklich Vergleichsschwellen abbildet, ohne aneurysma- oder
patientenspezifische Zielwerte zu suggerieren.
\end{abstract}

\begin{IEEEkeywords}
Blutdruck, HBPM, Zeitreihen, Aggregation, Median, Interquartilsabstand,
LaTeX, PGFPlots, Reproduzierbarkeit.
\end{IEEEkeywords}

\section{Einleitung und Motivation}
Die Selbstmessung des Blutdrucks zu Hause liefert über Wochen oder Monate eine
unregelmäßig getaktete Zeitreihe. Zwei methodische Probleme treten dabei
regelmäßig auf. Erstens variiert die Anzahl der Messungen pro Tag stark; ein
Tag mit acht Messungen würde in einer naiven Gesamtmittelung achtmal so stark
gewichtet wie ein Tag mit einer einzigen Messung. Zweitens enthalten reale
HBPM-Daten einzelne stark abweichende Werte (etwa unmittelbar nach körperlicher
Belastung), die robuste und nicht-robuste Lagemaße unterschiedlich beeinflussen.

\texttt{generate\_bp\_tikz.py} adressiert beide Punkte durch eine
deterministische, parametrisierbare Pipeline. Die Eingabe ist eine CSV-Datei,
die Ausgabe besteht aus einem in ein bestehendes LaTeX-Dokument einbindbaren
Fragment, zwei eigenständigen Standalone-Dokumenten sowie zwei optionalen
Statistik-CSVs. Sämtliche Lage- und Streuungsmaße werden im Quelltext explizit
und nachvollziehbar berechnet, sodass die Abbildungen unmittelbar gegen die
zugrunde liegenden Zahlen geprüft werden können.

\section{Datenmodell und Einlesen}
Erwartet wird eine CSV-Datei mit mindestens den Spalten \emph{Datum},
\emph{Systolisch} und \emph{Diastolisch}; optionale Spalten \emph{Puls},
\emph{Zeit} und \emph{Notiz} werden erkannt und, wo sinnvoll, ausgewertet. Die
Spaltenerkennung erfolgt über Aliasnamen in Deutsch und Englisch (etwa
\texttt{date}/\texttt{datum}, \texttt{systolic}/\texttt{systolisch}), wobei
Groß/Kleinschreibung und Sonderzeichen normalisiert werden.

Datumsangaben werden in den Formaten \texttt{YYYY-MM-DD}, \texttt{DD.MM.YYYY},
\texttt{DD.MM.YY}, \texttt{DD/MM/YYYY} und \texttt{DD/MM/YY} akzeptiert.
Zahlenwerte werden sowohl in englischer als auch in europäischer Schreibweise
geparst: Werte wie \texttt{130}, \texttt{130.5}, \texttt{130,5},
\texttt{1.234,56} und \texttt{1,234.56} werden korrekt interpretiert, indem das
zuletzt auftretende Trennzeichen als Dezimaltrenner gedeutet wird. Nicht
numerische Zusätze wie ``mmHg'' werden toleriert. Das Trennzeichen der
Eingabedatei wird standardmäßig automatisch erkannt (Komma, Semikolon oder
Tabulator) und kann mit \param{-{}-delimiter} erzwungen werden.

Beim Einlesen werden ausschließlich Zeilen im gewählten Datumsfenster
(\param{-{}-date-from} bis optional \param{-{}-date-to}) übernommen. Fehlende
oder ungültige systolische bzw.\ diastolische Werte führen zu einem klaren
Fehlerabbruch mit Zeilennummer, statt stillschweigend verfälschte Statistiken
zu erzeugen.

\section{Methodik der Aggregation}

\subsection{Stufe 1: Aggregation pro Kalendertag}
Im ersten Schritt werden alle Messungen eines Kalendertags zu einer
Tagesstatistik zusammengefasst. Für systolischen und diastolischen Druck sowie
für den abgeleiteten Pulsdruck (Differenz aus systolisch und diastolisch)
werden Minimum, Maximum, Median und arithmetisches Mittel bestimmt; für den
Puls analog, sofern vorhanden. Diese Tagesreduktion ist der entscheidende
Schritt gegen die Übergewichtung messintensiver Tage: Unabhängig davon, ob ein
Tag eine oder acht Messungen enthält, geht er anschließend als \emph{ein}
Tageswert in die weiteren Berechnungen ein.

Das Tagesdiagramm (Abbildung~1 der Ausgabe) zeigt je Tag den Tagesmedian als
Punkt und die Tagesspannweite (Minimum bis Maximum) als vertikalen Balken. Tage
mit nur einer Messung besitzen daher keine sichtbare Spannweite. Der Median
wird hier dem Mittelwert vorgezogen, weil er innerhalb eines Tages robust gegen
einen einzelnen Ausreißer ist und gemeinsam mit der dargestellten Spannweite
sowohl die typische Lage als auch die Streuung des Tages abbildet.

\subsection{Stufe 2: Blockaggregation}
Im zweiten Schritt werden die Tageswerte in Blöcke fester Länge gruppiert
(Standard sieben Tage, einstellbar über \param{-{}-block-days}). Die
Blockgrenzen werden ab \param{-{}-date-from} gezählt, sodass der erste Block
unabhängig vom Wochentag beginnt. Der letzte Block kann kürzer ausfallen; die
Bildunterschrift weist explizit darauf hin, wenn er weniger Tage umfasst und
daher weniger belastbar ist.

Für jeden Block werden zwei unterschiedliche Größen berechnet, die im Diagramm
bewusst getrennt dargestellt werden:
\begin{itemize}
\item die \emph{Zentrallinie} als Lagemaß des Blocks und
\item der \emph{Interquartilsabstand (IQR)} der Tagesmediane als Streuungsmaß,
dargestellt als Fehlerbalken von $Q_1$ bis $Q_3$.
\end{itemize}
Die Quartile werden über lineare Interpolation der sortierten Tagesmediane
bestimmt, vergleichbar mit der Voreinstellung gängiger Statistikpakete.

\section{Berechnungsvarianten der Zentrallinie}
Die zentrale methodische Entscheidung betrifft das Lagemaß der Blockzentrallinie.
Das Werkzeug bietet über \param{-{}-week-central} zwei Varianten an, die jeweils
einen klaren Zweck verfolgen.

\subsection{Mittelwert (Voreinstellung)}
In der Voreinstellung (\texttt{mean}) ist die Zentrallinie der nach
Kalendertagen gewichtete Mittelwert der Tagesmittelwerte. Formal gilt für einen
Block mit $D$ Tagen und Tagesmittelwerten $\bar{x}_d$:
\begin{equation}
\mu_{\mathrm{Block}} = \frac{1}{D}\sum_{d=1}^{D} \bar{x}_d .
\end{equation}
Da bereits $\bar{x}_d$ das Mittel der Messungen des Tages $d$ ist, ist
$\mu_{\mathrm{Block}}$ ein zweistufiges, nach Tagen gleichgewichtetes Mittel.
Diese Variante ist anschlussfähig an die klinisch gebräuchlichen
HBPM/ESC-Mittelwertschwellen, die ebenfalls auf Mittelwerten beruhen, und
eignet sich daher für den direkten Vergleich mit solchen Schwellen.

\subsection{Median der Tagesmediane}
Mit \texttt{median} ist die Zentrallinie der Median der Tagesmediane des Blocks:
\begin{equation}
m_{\mathrm{Block}} = \operatorname{median}\bigl(\{\tilde{x}_d\}_{d=1}^{D}\bigr),
\end{equation}
wobei $\tilde{x}_d$ der Median des Tages $d$ ist. Diese Variante ist robuster
gegen einzelne Ausreißertage und ist zudem konsistent zur dargestellten
IQR-Box, da beide auf denselben Tagesmedianen beruhen. Sie empfiehlt sich,
wenn die typische Lage trotz vereinzelter Belastungsspitzen abgebildet werden
soll.

\subsection{Empfehlung}
Die Wahl ist keine Frage von richtig oder falsch, sondern von Zweck: Für den
Abgleich mit mittelwertbasierten Schwellen ist \texttt{mean} naheliegend, für
eine ausreißerrobuste Verlaufsdarstellung \texttt{median}. Da Streuung
(IQR-Box) und Lage getrennt dargestellt werden, bleibt die Information in
beiden Fällen vollständig sichtbar.

\section{Ausreißer-Markierung}
Optional (\param{-{}-week-outliers}) markiert das Wochendiagramm einzelne Tage,
deren Tagesmedian eine klinische Schwelle über- oder unterschreitet, als kleine
Punkte neben dem jeweiligen IQR-Marker. Die oberen Schwellen sind
voreingestellt auf 135\,mmHg systolisch und 85\,mmHg diastolisch und
entsprechen damit den im Diagramm eingezeichneten HBPM-Vergleichslinien;
sie sind über \param{-{}-week-outlier-sys-hi} und
\param{-{}-week-outlier-dia-hi} anpassbar. Zusätzlich lassen sich über
\param{-{}-week-outlier-sys-lo} und \param{-{}-week-outlier-dia-lo} optionale
\emph{untere} Schwellen aktivieren, um etwa unter Therapie zu niedrige Tage
sichtbar zu machen. Diese Markierung ersetzt keine statistische
Ausreißerdefinition, sondern hebt klinisch interessante Einzeltage hervor,
ohne sie aus der Aggregation zu entfernen; alle Tage bleiben in Mittelwert,
Median und IQR enthalten.

\section{Referenzelemente und Schwellen}
Beide Diagramme enthalten dieselben Referenzelemente, die in einem gemeinsamen
Absatz einmal erläutert und von den Bildunterschriften referenziert werden. Die
hellgrauen Korridore 120--129\,mmHg systolisch und 70--79\,mmHg diastolisch sind
allgemeine ESC-Orientierungen unter Therapie bei individueller Verträglichkeit
und ausdrücklich \emph{keine} aneurysmaspezifischen Zielwertlinien. Die
punktierten Linien bei 135\,mmHg systolisch und 85\,mmHg diastolisch markieren
häufig verwendete häusliche Vergleichsschwellen. Diese bewusste Trennung
zwischen Vergleichsschwelle und individuellem Therapieziel verhindert, dass die
Grafik eine nicht belegte, patientenspezifische Zielwertlinie suggeriert.

Die automatische Zusammenfassungszeile am unteren Rand des Tagesdiagramms nennt
den nach Tagen gewichteten Mittelwert, den Median der Tagesmediane und die
Anzahl der Tage mit Tagesmittel unter 135/85\,mmHg. Sie kann mit
\param{-{}-no-daily-summary-label} unterdrückt werden.

\section{Erzeugte Dateien}
Standardmäßig schreibt das Werkzeug ein einbindbares LaTeX-Fragment
(\param{-{}-out}), ein einseitiges Standalone-Dokument mit beiden Diagrammen
untereinander (\param{-{}-standalone-out}) sowie ein ``zweiseitiges''
Standalone-Dokument, in dem jedes Diagramm auf einer eng beschnittenen eigenen
Seite steht (\param{-{}-two-sides-out}). Optional werden zwei Statistik-CSVs
geschrieben (\param{-{}-daily-stats}, \param{-{}-weekly-stats}). Das einseitige
Dokument skaliert die Diagrammhöhen zur Laufzeit aus der verfügbaren Seitenhöhe,
sodass die Seite unabhängig von der Länge der Bildunterschriften vollständig
ausgenutzt wird.

Tabelle~\ref{tab:daily_cols} und Tabelle~\ref{tab:weekly_cols} dokumentieren die
Spalten der beiden Statistik-CSVs.

\begin{table}[!t]
\caption{Spalten der Tagesstatistik-CSV (\texttt{-{}-daily-stats}).}
\label{tab:daily_cols}
\centering
\footnotesize
\renewcommand{\arraystretch}{1.15}
\begin{tabular}{>{\raggedright\arraybackslash}p{0.30\columnwidth} >{\raggedright\arraybackslash}p{0.58\columnwidth}}
\toprule
\textbf{Spalte} & \textbf{Bedeutung} \\
\midrule
\texttt{date} & Kalendertag (ISO 8601) \\
\texttt{day\_index} & Tagesindex ab \texttt{date-from} (0-basiert) \\
\texttt{n} & Anzahl Messungen an diesem Tag \\
\texttt{sys\_min/max} & Minimum/Maximum systolisch \\
\texttt{sys\_median} & Tagesmedian systolisch \\
\texttt{sys\_mean} & Tagesmittel systolisch \\
\texttt{dia\_min/max} & Minimum/Maximum diastolisch \\
\texttt{dia\_median} & Tagesmedian diastolisch \\
\texttt{dia\_mean} & Tagesmittel diastolisch \\
\texttt{pulse\_*} & Puls min/max/median/mean (falls vorhanden) \\
\texttt{pulse\_pressure\_*} & Pulsdruck min/max/median/mean \\
\bottomrule
\end{tabular}
\end{table}

\begin{table}[!t]
\caption{Spalten der Blockstatistik-CSV (\texttt{-{}-weekly-stats}).}
\label{tab:weekly_cols}
\centering
\footnotesize
\renewcommand{\arraystretch}{1.15}
\begin{tabular}{>{\raggedright\arraybackslash}p{0.40\columnwidth} >{\raggedright\arraybackslash}p{0.48\columnwidth}}
\toprule
\textbf{Spalte} & \textbf{Bedeutung} \\
\midrule
\texttt{block} & Blockindex (0-basiert) \\
\texttt{start\_date}, \texttt{end\_date} & Anfang/Ende des Blocks \\
\texttt{x\_mid} & x-Position (Tagesindex) der Blockmitte \\
\texttt{n\_days} & Anzahl Tage im Block \\
\texttt{n\_readings} & Anzahl Messungen im Block \\
\texttt{label} & Achsenbeschriftung des Blocks \\
\texttt{sys\_mean\_of\_daily\_means} & gewichteter Mittelwert systolisch \\
\texttt{sys\_q1/q3\_daily\_medians} & $Q_1$/$Q_3$ der Tagesmediane systolisch \\
\texttt{dia\_mean\_of\_daily\_means} & gewichteter Mittelwert diastolisch \\
\texttt{dia\_q1/q3\_daily\_medians} & $Q_1$/$Q_3$ der Tagesmediane diastolisch \\
\texttt{pulse\_pressure\_*} & analog für den Pulsdruck \\
\bottomrule
\end{tabular}
\end{table}

\section{Anonymisiertes Anwendungsbeispiel}
Die Abbildungen~\ref{fig:daily} und~\ref{fig:weekly} zeigen die beiden
erzeugten Diagramme an einem anonymisierten, realitätsnahen Datensatz
(101 Messungen an 37 Kalendertagen). Abbildung~\ref{fig:daily} ist der
Tagesverlauf mit Tagesmedian, Tagesspannweite und der automatischen
Zusammenfassungszeile sowie den \texttt{n=}-Labels.
Abbildung~\ref{fig:weekly} ist die verdichtete 7-Tage-Übersicht; hier wurde die
Median-Zentrallinie (\param{-{}-week-central median}) und die
Ausreißermarkierung (\param{-{}-week-outliers}) gewählt. In beiden Abbildungen
sind die ESC-Korridore und HBPM-Vergleichslinien als gemeinsame Referenz
eingezeichnet.

\begin{figure*}[!t]
\centering
\includegraphics[width=0.92\textwidth]{fig1.pdf}
\caption{Anonymisiertes Beispiel, Tagesverlauf: Punkte zeigen den Tagesmedian,
vertikale Balken die Tagesspannweite (Min--Max); graue \texttt{n=}-Angaben oben
die Messungen pro Tag. Die Zusammenfassungszeile nennt gewichteten Mittelwert,
Median der Tagesmediane und die Anzahl Tage unter 135/85\,mmHg.}
\label{fig:daily}
\end{figure*}

\begin{figure*}[!t]
\centering
\includegraphics[width=0.92\textwidth]{fig2.pdf}
\caption{Anonymisiertes Beispiel, 7-Tage-Übersicht mit Median-Zentrallinie und
Ausreißermarkierung: Linien zeigen den Median der Tagesmediane, Fehlerbalken den
Interquartilsbereich der Tagesmediane; Kreise markieren einzelne Tage, deren
Tagesmedian die Schwelle (135/85\,mmHg) überschreitet.}
\label{fig:weekly}
\end{figure*}

Abbildung~\ref{fig:onepage} zeigt zusätzlich das einseitige Standalone-Dokument
(\param{-{}-standalone-out}), in dem beide Diagramme samt gemeinsamem
Referenzabsatz und Bildunterschriften auf einer Seite zusammengeführt sind. Die
Diagrammhöhen werden dabei zur Laufzeit aus der verfügbaren Seitenhöhe
berechnet, sodass die Seite unabhängig von der Länge der Bildunterschriften
vollständig ausgenutzt wird.

\begin{figure}[!t]
\centering
\fbox{\includegraphics[width=0.97\columnwidth]{fig3.pdf}}
\caption{Anonymisiertes Beispiel des einseitigen Standalone-Dokuments
(\texttt{-{}-standalone-out}): beide Diagramme, der gemeinsame Referenzabsatz
und die Bildunterschriften auf einer Seite, mit zur Seitenhöhe skalierten
Diagrammen.}
\label{fig:onepage}
\end{figure}

Schließlich zeigt Abbildung~\ref{fig:twosides} die zweiseitige Variante
(\param{-{}-two-sides-out}), bei der jedes Diagramm auf einer eigenen, eng
beschnittenen Seite steht und die zugehörige Bildunterschrift als Textknoten
unterhalb der Achse gesetzt wird. Dieses Format eignet sich besonders zum
Einbetten der einzelnen Diagramme als separate Anhangsabbildungen.

\begin{figure*}[!t]
\centering
\begin{subfigure}{0.48\textwidth}
\centering
\fbox{\includegraphics[width=\textwidth]{fig4a.pdf}}
\caption{Seite 1: Tagesverlauf.}
\label{fig:twosides_a}
\end{subfigure}
\hfill
\begin{subfigure}{0.48\textwidth}
\centering
\fbox{\includegraphics[width=\textwidth]{fig4b.pdf}}
\caption{Seite 2: 7-Tage-Übersicht.}
\label{fig:twosides_b}
\end{subfigure}
\caption{Anonymisiertes Beispiel der zweiseitigen Variante
(\texttt{-{}-two-sides-out}): jedes Diagramm steht auf einer eigenen,
beschnittenen Seite mit der Bildunterschrift als Knoten unterhalb der Achse.}
\label{fig:twosides}
\end{figure*}

\section{Parameterreferenz}
Tabelle~\ref{tab:params} listet alle Kommandozeilenparameter mit Voreinstellung
und Zweck.

\begin{table*}[!t]
\caption{Vollständige Parameterreferenz von \texttt{generate\_bp\_tikz.py}.}
\label{tab:params}
\centering
\footnotesize
\renewcommand{\arraystretch}{1.25}
\begin{tabular}{>{\raggedright\arraybackslash}p{0.24\textwidth} >{\raggedright\arraybackslash}p{0.16\textwidth} >{\raggedright\arraybackslash}p{0.52\textwidth}}
\toprule
\textbf{Parameter} & \textbf{Voreinstellung} & \textbf{Zweck} \\
\midrule
\param{-{}-csv} & --- (Pflicht) & Eingabe-CSV mit Datum, Systolisch, Diastolisch (optional Zeit/Puls/Notiz). \\
\param{-{}-date-from} & --- (Pflicht) & Erstes Datum im Auswertungsfenster. \\
\param{-{}-date-to} & alle ab \texttt{date-from} & Letztes Datum; bei Auslassung alle nachfolgenden Werte. \\
\param{-{}-out} & \texttt{bp\_diagrams.tex} & Pfad des einbindbaren LaTeX-Fragments. \\
\param{-{}-daily-stats} & aus & Optionale Tagesstatistik als CSV. \\
\param{-{}-weekly-stats} & aus & Optionale Blockstatistik als CSV. \\
\param{-{}-block-days} & \texttt{7} & Blocklänge der zweiten Aggregationsstufe in Tagen. \\
\param{-{}-delimiter} & \texttt{auto} & Trennzeichen: \texttt{auto}/\texttt{comma}/\texttt{semicolon}/\texttt{tab}. \\
\param{-{}-standalone-out} & \texttt{bp\_\hspace{0pt}diagrams\_\hspace{0pt}both\_\hspace{0pt}onepage\_\hspace{0pt}standalone.tex} & Einseitiges Standalone mit beiden Diagrammen. \\
\param{-{}-no-standalone} & aus & Einseitiges Standalone nicht schreiben. \\
\param{-{}-standalone-title} & \texttt{Blutdruckdiagramme} & Titel über den Standalone-Diagrammen. \\
\param{-{}-show-daily-n} & aus & Graue \texttt{n=}-Labels (Messungen pro Tag) im Tagesdiagramm. \\
\param{-{}-daily-n-y} & automatisch & Feste y-Position der \texttt{n=}-Labels. \\
\param{-{}-week-outliers} & aus & Ausreißertage im Wochendiagramm markieren. \\
\param{-{}-week-outlier-sys-hi} & \texttt{135.0} & Obere systolische Ausreißerschwelle (mmHg). \\
\param{-{}-week-outlier-dia-hi} & \texttt{85.0} & Obere diastolische Ausreißerschwelle (mmHg). \\
\param{-{}-week-outlier-sys-lo} & aus & Optionale untere systolische Schwelle. \\
\param{-{}-week-outlier-dia-lo} & aus & Optionale untere diastolische Schwelle. \\
\param{-{}-week-central} & \texttt{mean} & Zentrallinie: \texttt{mean} (gewichteter Mittelwert) oder \texttt{median} (Median der Tagesmediane). \\
\param{-{}-no-daily-summary-label} & aus & Automatische Zusammenfassungszeile im Tagesdiagramm unterdrücken. \\
\param{-{}-two-sides-out} & \texttt{bp\_\hspace{0pt}diagrams\_\hspace{0pt}standalone\_\hspace{0pt}two\_\hspace{0pt}sides.tex} & Zweiseitiges Standalone (je Diagramm eine beschnittene Seite). \\
\param{-{}-no-two-sides} & aus & Zweiseitiges Standalone nicht schreiben. \\
\param{-{}-two-sides-width-cm} & \texttt{16.0} & Feste Achsenbreite (cm) der zweiseitigen Variante. \\
\bottomrule
\end{tabular}
\end{table*}

\section{Bedienung}
\subsection{Voraussetzungen}
Das einbindende LaTeX-Dokument benötigt:
\begin{lstlisting}[style=sh]
\usepackage{tikz}
\usepackage{pgfplots}
\pgfplotsset{compat=1.18}
\end{lstlisting}
Das Skript selbst hat keine externen Python-Abhängigkeiten über die
Standardbibliothek hinaus.

\subsection{Minimaler Aufruf}
\begin{lstlisting}[style=sh]
python3 generate_bp_tikz.py \
  --csv bp_raw.csv \
  --date-from 2026-05-15
\end{lstlisting}
Erzeugt Fragment und beide Standalone-Dokumente mit den Standardnamen.

\subsection{Mit Datumsfenster und Statistik-Export}
\begin{lstlisting}[style=sh]
python3 generate_bp_tikz.py \
  --csv bp_raw.csv \
  --date-from 15.05.2026 \
  --date-to 20.06.2026 \
  --out bp_fragment.tex \
  --daily-stats bp_daily_stats.csv \
  --weekly-stats bp_weekly_stats.csv
\end{lstlisting}

\subsection{Median-Zentrallinie mit Ausreißern}
\begin{lstlisting}[style=sh]
python3 generate_bp_tikz.py \
  --csv bp_raw.csv \
  --date-from 2026-05-15 \
  --week-central median \
  --week-outliers
\end{lstlisting}

\subsection{Voll bestückter Aufruf}
\begin{lstlisting}[style=sh]
python3 generate_bp_tikz.py \
  --csv bp_raw.csv \
  --date-from 2026-05-15 \
  --date-to 2026-06-20 \
  --week-central median --week-outliers \
  --week-outlier-sys-hi 135 \
  --week-outlier-dia-hi 85 \
  --show-daily-n \
  --out bp_fragment.tex \
  --standalone-out bp_onepage.tex \
  --two-sides-out bp_two_sides.tex \
  --two-sides-width-cm 16 \
  --daily-stats bp_daily_stats.csv \
  --weekly-stats bp_weekly_stats.csv
\end{lstlisting}

\subsection{Einbindung der Ausgabe}
Das Fragment wird mit \verb|\input{bp_fragment.tex}| in den Fließtext
eingebunden. Alternativ kann das einseitige oder zweiseitige Standalone separat
übersetzt und das resultierende PDF als Anhangsabbildung eingebettet werden.

\section{Reproduzierbarkeit und Validierung}
Da sämtliche Lage- und Streuungsmaße im Quelltext explizit berechnet und über
die Statistik-CSVs ausgegeben werden, lassen sich die Diagramme unmittelbar
gegen die Zahlen prüfen. Ein typischer Validierungsschritt besteht darin, die
Zusammenfassungszeile des Tagesdiagramms (gewichteter Mittelwert, Median der
Tagesmediane, Anzahl Tage unter 135/85\,mmHg) mit den Werten in der
Tagesstatistik-CSV abzugleichen. Gleiche Eingabe, gleiches Datumsfenster und
gleiche Parameter führen deterministisch zu identischer Ausgabe.

\section{Implementierung und getestete Umgebung}
Das Werkzeug ist in Python~3 implementiert und verwendet ausschließlich die
Standardbibliothek (keine Drittpaket-Abhängigkeiten); empfohlen wird
Python~3.9 oder neuer. Die erzeugten Dokumente wurden mit \texttt{pdflatex}
(TeX~Live) unter Verwendung von \texttt{tikz} und \texttt{pgfplots}
(\texttt{compat=1.18}) gesetzt; die beigefügten IEEE-Papers nutzen
\texttt{IEEEtran.cls} in Version~V1.8b.

Getestet wurde ausschließlich auf folgenden Plattformen: (i)~macOS mit
Desktop-Python und einer Desktop-TeX-Live-/MacTeX-Installation; (ii)~iPad und
iPhone mit \emph{a-Shell}, einem lokalen iOS-Terminal, das Python~3 und
TeX~Live (mit TikZ und LuaTeX) bündelt -- sowohl das Skript als auch die
LaTeX-Übersetzung laufen dort lokal auf dem Gerät; sowie (iii)~\emph{Texifier}
(vormals TeXpad) auf macOS und iPadOS/iOS zum Bearbeiten und Setzen der
LaTeX-Dokumente. Zu beachten ist, dass der integrierte Live-Setzer von Texifier
(TexpadTeX) in einer Sandbox läuft und keinen Shell-Escape (\texttt{\textbackslash write18})
erlaubt; das Python-Skript selbst wird daher in a-Shell bzw.\ in einer
gewöhnlichen Shell ausgeführt, während Texifier zum Setzen der erzeugten
\texttt{.tex}-Dateien dient. Andere Plattformen (Windows, Linux) sollten mit
einer üblichen Python-3- und TeX-Live-Installation funktionieren, wurden aber
nicht getestet.

\section{Einschränkungen}
Das Werkzeug ist ein Aufbereitungs- und Darstellungswerkzeug, kein
diagnostisches System. Die eingezeichneten Schwellen sind Vergleichs- und
Orientierungslinien; die Festlegung individueller Zielwerte obliegt der
behandelnden Ärztin oder dem behandelnden Arzt. Die Aussagekraft kurzer Blöcke
mit wenigen Tagen ist begrenzt und wird in der Bildunterschrift entsprechend
gekennzeichnet. Die Ausreißermarkierung ist eine schwellenbasierte
Hervorhebung und keine statistische Ausreißerklassifikation.

\section{Zusammenfassung}
\texttt{generate\_bp\_tikz.py} überführt unregelmäßig getaktete
HBPM-Zeitreihen über eine zweistufige Aggregation in zwei abgestimmte,
reproduzierbare Diagramme. Die getrennte Darstellung von Lage (Mittelwert oder
Median) und Streuung (Spannweite bzw.\ IQR), die transparente Ausreißerlogik
und die explizite Trennung von Vergleichsschwelle und Therapieziel erlauben
eine klinisch interpretierbare und zugleich nachvollziehbare Darstellung. Die
vollständige Parametrisierung über die Kommandozeile macht das Werkzeug
sowohl für einmalige Auswertungen als auch für wiederkehrende, automatisierte
Berichte geeignet.

\section*{Danksagung und Offenlegung zur Nutzung generativer KI}
Im Sinne der IEEE-Richtlinien zur Nutzung KI-generierter Inhalte wird Folgendes
offengelegt. Zur Erstellung dieses Dokuments wurde das KI-System Claude
(Anthropic) eingesetzt. Die KI wurde unterstützend verwendet bei: der
sprachlichen Ausarbeitung und Strukturierung des Fließtextes, der Formulierung
der Methodikbeschreibung auf Grundlage des vom Autor entwickelten Quelltextes,
der Erstellung der LaTeX- und Tabellenstruktur sowie der Aufbereitung der
anonymisierten Beispielabbildungen. Konzeption, Software, zugrunde liegende
Daten, fachliche Aussagen und die endgültige inhaltliche Verantwortung liegen
vollständig beim Autor; sämtliche KI-gestützten Beiträge wurden vom Autor
geprüft, verifiziert und gegebenenfalls korrigiert. Die KI ist kein Autor dieser
Arbeit. Eine ausführlichere Beschreibung des KI-Anteils findet sich in der Datei
\texttt{AI\_USAGE.md} des zugehörigen Repositories.

\end{document}
